\section{测地线的最短性}\label{sec:minimizing-properties-of-geodesics}

设 $M$ 是带有仿射联络 $\nabla$ 的可微流形. 如第~\ref{sec:affine-connections}~节所见, 给定点 $p\in M$ 和切向量 $v\in T_pM$, 存在唯一的测地线 $c_v:I\to M$, 定义在极大开区间 $I\subset\mathbb{R}$ 上, 使得 $0\in I$, $c_v(0)=p$ 且 $\dot c_v(0)=v$. 现考虑由 $\gamma(t)=c_v(at)$ 定义的曲线 $\gamma:J\to M$, 其中 $a\in\mathbb{R}$, 而 $J$ 是映射 $t\mapsto at$ 下 $I$ 的原像. 有
\[
  \dot\gamma(t)=a\dot c_v(at),
\]
从而
\[
  \nabla_{\dot\gamma}\dot\gamma
  =\nabla_{a\dot c_v}(a\dot c_v)=a^2\nabla_{\dot c_v}\dot c_v=0.
\]
所以 $\gamma$ 也是测地线. 由于 $\gamma(0)=c_v(0)=p$ 且 $\dot\gamma(0)=a\dot c_v(0)=av$, $\gamma$ 就是初速度为 $av\in T_pM$ 的唯一测地线 (即 $\gamma=c_{av}$). 因此对所有 $t\in I$, 有 $c_{av}(t)=c_v(at)$. 这一性质有时称为\textbf{测地线的齐次性}. 注意, 取 $a$ 足够小, 就可以使区间 $J$ 任意大. 若 $1\in I$, 定义 $\exp_p(v)=c_v(1)$. 由测地线的齐次性, 可以对 $T_pM$ 中原点的某个开邻域 $U$ 内的 $v$ 定义 $\exp_p(v)$. 由此得到的映射 $\exp_p:U\subset T_pM\to M$ 称为 $p$ 处的\textbf{指数映射} (Fig.~\ref{fig:exponential-map}).
\RGfigure[.53\textwidth]{exponential-map}{指数映射}

\begin{RiemannProposition}\label{prop:exponential-map-local-diffeomorphism}
存在包含原点的开集 $U\subset T_pM$, 使 $\exp_p:U\to M$ 是到某个包含 $p$ 的开集 $V\subset M$ 的微分同胚 ($V$ 称为\textbf{正规邻域}).
\end{RiemannProposition}

\begin{RiemannProof}
由常微分方程的解对初始数据的光滑依赖性 (参见 \RGcite{ch03-Arn92}), 指数映射显然可微. 若 $v\in T_pM$ 使 $\exp_p(v)$ 有定义, 则由齐次性有 $\exp_p(tv)=c_{tv}(1)=c_v(t)$. 因而
\[
  (\mathrm{d}\exp_p)_0v=\left.\frac{\mathrm{d}}{\mathrm{d}t}\exp_p(tv)\right|_{t=0}
  =\left.\frac{\mathrm{d}}{\mathrm{d}t}c_v(t)\right|_{t=0}=v.
\]
因此 $(\mathrm{d}\exp_p)_0:T_0(T_pM)\cong T_pM\to T_pM$ 是恒等映射. 由反函数定理, $\exp_p$ 是从 $0\in T_pM$ 的某个开邻域 $U$ 到包含 $p=\exp_p(0)$ 的某个开集 $V\subset M$ 的微分同胚.
\end{RiemannProof}

\begin{RiemannExample}\label{ex:exponential-map-of-standard-sphere}
考虑带标准度量的 $S^2$ 上的 Levi-Civita 联络, 设 $p\in S^2$. 则 $\exp_p(v)$ 对所有 $v\in T_pS^2$ 都有定义, 但它不是微分同胚, 因为显然它不是单射. 然而, 它在开球 $B_\pi(0)\subset T_pS^2$ 上的限制是到 $S^2\setminus\{-p\}$ 的微分同胚.
\end{RiemannExample}

现在设 $(M,\langle\cdot,\cdot\rangle)$ 是 Riemann 流形, $\nabla$ 是其 Levi-Civita 联络. 由于 $\langle\cdot,\cdot\rangle$ 在 $T_pM$ 上定义了一个内积, 可以把 $T_pM$ 看作 Euclidean $n$ 维空间 $\mathbb{R}^n$. 设 $E$ 是在 $T_pM\setminus\{0\}$ 上由下式定义的向量场:
\[
  E_v=\frac{v}{\lVert v\rVert},
\]
并在 $V\setminus\{p\}$ 上定义 $X:=(\exp_p)_*E$, 其中 $V\subset M$ 是一个正规邻域. 有
\begin{align*}
  X_{\exp_p(v)}&=(\mathrm{d}\exp_p)_vE_v
   =\left.\frac{\mathrm{d}}{\mathrm{d}t}\exp_p\left(v+t\frac{v}{\lVert v\rVert}\right)\right|_{t=0}\\
  &=\left.\frac{\mathrm{d}}{\mathrm{d}t}c_v\left(1+\frac{t}{\lVert v\rVert}\right)\right|_{t=0}
   =\frac{1}{\lVert v\rVert}\dot c_v(1).
\end{align*}
由于 $\lVert\dot c_v(1)\rVert=\lVert\dot c_v(0)\rVert=\lVert v\rVert$, 可知 $X_{\exp_p(v)}$ 是测地线 $c_v$ 的单位切向量. 特别地, $X$ 必须满足
\[
  \nabla_XX=0.
\]
对使 $\overline{B_\varepsilon(0)}\subset U:=\exp_p^{-1}(V)$ 的 $\varepsilon>0$, 将以 $p$ 为中心, 半径为 $\varepsilon$ 的\phantomsection\label{idx:normal-balls}\textbf{正规球}定义为开集 $B_\varepsilon(p):=\exp_p(B_\varepsilon(0))$, 将以 $p$ 为中心, 半径为 $\varepsilon$ 的\textbf{正规球面}定义为紧子流形 $S_\varepsilon(p):=\exp_p(\partial B_\varepsilon(0))$. 现在证明 $X$ (从而经过 $p$ 的测地线) 与正规球面正交. 为此, 选取局部参数化 $\varphi:W\subset\mathbb{R}^{n-1}\to S^{n-1}\subset T_pM$, 并借此定义参数化 $\widetilde\varphi:(0,+\infty)\times W\to T_pM$:
\[
  \widetilde\varphi(r,\theta^1,\ldots,\theta^{n-1})=r\varphi(\theta^1,\ldots,\theta^{n-1})
\]
(从而 $(r,\theta^1,\ldots,\theta^{n-1})$ 是 $T_pM$ 上的球坐标). 注意
\[
  \frac{\partial}{\partial r}=E,
\]
因为
\[
  E_{\widetilde\varphi(r,\theta)}=E_{r\varphi(\theta)}
  =\varphi(\theta)=\frac{\partial\widetilde\varphi}{\partial r}(r,\theta),
\]
所以
\begin{equation}\label{eq:radial-field-in-normal-ball}
  X=(\exp_p)_*\frac{\partial}{\partial r}.
\end{equation}
由于 $\frac{\partial}{\partial\theta^i}$ 与 $\{r=\varepsilon\}$ 相切, 向量场
\begin{equation}\label{eq:angular-field-in-normal-ball}
  Y_i:=(\exp_p)_*\frac{\partial}{\partial\theta^i}
\end{equation}
与 $S_\varepsilon(p)$ 相切. 还注意到, $\bigl\lVert\frac{\partial}{\partial\theta^i}\bigr\rVert
=\bigl\lVert\frac{\partial\widetilde\varphi}{\partial\theta^i}\bigr\rVert
=r\bigl\lVert\frac{\partial\varphi}{\partial\theta^i}\bigr\rVert$ 与 $r$ 成正比. 因而当 $r\to0$ 时, $\frac{\partial}{\partial\theta^i}\to0$, 这就意味着当 $q\to p$ 时, $(Y_i)_q\to0_p$. 由于 $\exp_p$ 是局部微分同胚, 向量场 $X$ 和 $Y_i$ 在每一点都线性无关. 此外,
\[
  [X,Y_i]=\left[(\exp_p)_*\frac{\partial}{\partial r},(\exp_p)_*\frac{\partial}{\partial\theta^i}\right]
  =(\exp_p)_*\left[\frac{\partial}{\partial r},\frac{\partial}{\partial\theta^i}\right]=0
\]
[参见第~\ref{chap:differentiable-manifolds}~章的 Exercise~\ref{ex:vector-fields}(\ref{ex:pushforward-preserves-lie-bracket})]. 又因 Levi-Civita 联络对称, 等价地有
\[
  \nabla_XY_i=\nabla_{Y_i}X.
\]
为证明 $X$ 与正规球面 $S_\varepsilon(p)$ 正交, 只需证明 $X$ 与每个向量场 $Y_i$ 正交. 事实上, 由于 $\nabla_XX=0$ 和 $\lVert X\rVert=1$, 有
\begin{align*}
  X\cdot\langle X,Y_i\rangle
  &=\langle\nabla_XX,Y_i\rangle+\langle X,\nabla_XY_i\rangle\\
  &=\langle X,\nabla_{Y_i}X\rangle=\frac12Y_i\cdot\langle X,X\rangle=0,
\end{align*}
因而 $\langle X,Y_i\rangle$ 沿每条经过 $p$ 的测地线为常数. 所以
\[
  \langle X,Y_i\rangle(\exp_p v)
  =\bigl\langle X_{\exp_p(v)},(Y_i)_{\exp_p(v)}\bigr\rangle
  =\lim_{t\to0}\bigl\langle X_{\exp_p(tv)},(Y_i)_{\exp_p(tv)}\bigr\rangle=0
\]
(因为 $\lVert X\rVert=1$, 且当 $q\to p$ 时, $(Y_i)_q\to0_p$). 所以经过 $p$ 的每条测地线都与所有以 $p$ 为中心的正规球面正交. 利用这一点, 得到以下结果.

\begin{RiemannProposition}\label{prop:length-bound-from-normal-sphere}
设 $\gamma:[0,1]\to M$ 是可微曲线, 满足 $\gamma(0)=p$ 和 $\gamma(1)\in S_\varepsilon(p)$, 其中 $S_\varepsilon(p)$ 是正规球面. 则 $l(\gamma)\geq\varepsilon$, 且 $l(\gamma)=\varepsilon$ 当且仅当 $\gamma$ 是一条重参数化的测地线.
\end{RiemannProposition}

\begin{RiemannProof}
可以假设对所有 $t\in(0,1)$ 均有 $\gamma(t)\ne p$, 否则很容易构造曲线 $\widetilde\gamma:[0,1]\to M$, 满足 $\widetilde\gamma(0)=p$, $\widetilde\gamma(1)=\gamma(1)\in S_\varepsilon(p)$ 和 $l(\widetilde\gamma)<l(\gamma)$. 同理, 可假设 $\gamma([0,1))\subset B_\varepsilon(p)$. 于是可以写成
\[
  \gamma(t):=\exp_p(r(t)n(t)),
\]
其中对 $t\in(0,1]$, $r(t)\in(0,\varepsilon]$ 和 $n(t)\in S^{n-1}$ 都有定义. 注意, $r(t)$ 可以延拓为 $[0,1]$ 上的光滑函数. 于是
\[
  \dot\gamma(t)=(\exp_p)_*(\dot r(t)n(t)+r(t)\dot n(t)).
\]
由于 $\langle n(t),n(t)\rangle=1$, 有 $\langle\dot n(t),n(t)\rangle=0$, 所以 $\dot n(t)$ 与 $\partial B_{r(t)}(0)$ 相切. 注意到 $n(t)=\bigl(\frac{\partial}{\partial r}\bigr)_{r(t)n(t)}$, 得到
\[
  \dot\gamma(t)=\dot r(t)X_{\gamma(t)}+Y(t),
\]
其中 $X=(\exp_p)_*\frac{\partial}{\partial r}$, 而 $Y(t)=r(t)(\exp_p)_*\dot n(t)$ 与 $S_{r(t)}(p)$ 相切, 因而与 $X_{\gamma(t)}$ 正交. 所以
\begin{align*}
  l(\gamma)
  &=\int_0^1\bigl\langle\dot r(t)X_{\gamma(t)}+Y(t),\dot r(t)X_{\gamma(t)}+Y(t)\bigr\rangle^{\frac12}\,\mathrm{d}t\\
  &=\int_0^1\left(\dot r(t)^2+\lVert Y(t)\rVert^2\right)^{\frac12}\,\mathrm{d}t\\
  &\geq\int_0^1\dot r(t)\,\mathrm{d}t=r(1)-r(0)=\varepsilon.
\end{align*}
显然, $l(\gamma)=\varepsilon$ 当且仅当对所有 $t\in[0,1]$ 有 $\lVert Y(t)\rVert\equiv0$ 和 $\dot r(t)\geq0$. 在这种情况下, $\dot n(t)=0$ (从而 $n$ 为常量), 而 $\gamma(t)=\exp_p(r(t)n)=c_{r(t)n}(1)=c_n(r(t))$ 在相差重参数化的意义下, 是通过 $p$ 且初始条件为 $n\in T_pM$ 的测地线.
\end{RiemannProof}

\begin{RiemannDefinition}\label{def:piecewise-differentiable-curve}
\textbf{分段可微曲线}是一个连续映射 $c:[a,b]\to M$, 存在 $a=t_0<t_1<\cdots<t_{n-1}<t_n=b$, 使得对 $i=1,\ldots,n$, $c$ 在 $[t_{i-1},t_i]$ 上的限制可微. 若 $c(a)=p$ 且 $c(b)=q$, 则称 $c$ \textbf{连接} $p\in M$ 和 $q\in M$.
\end{RiemannDefinition}

分段可微曲线的\textbf{长度}没有定义上的困难. 显然, 若现在允许分段可微的重参数化, Proposition~\ref{prop:length-bound-from-normal-sphere} 很容易推广到分段可微曲线. 利用 Proposition~\ref{prop:length-bound-from-normal-sphere} 的这一推广, 指数映射的性质以及长度在重参数化下的不变性, 容易证明以下结果.

\begin{RiemannTheorem}\label{thm:radial-geodesic-minimizes-length}
设 $(M,\langle\cdot,\cdot\rangle)$ 是 Riemann 流形, $p\in M$, 且 $B_\varepsilon(p)$ 是以 $p$ 为中心的正规球. 则对每一点 $q\in B_\varepsilon(p)$, 存在连接 $p$ 到 $q$ 的测地线 $c:I\to M$. 此外, 若 $\gamma:J\to M$ 是连接 $p$ 到 $q$ 的任意其他分段可微曲线, 则 $l(\gamma)\geq l(c)$, 且 $l(\gamma)=l(c)$ 当且仅当 $\gamma$ 是 $c$ 的重参数化.\hfill\openbox
\end{RiemannTheorem}

反过来, 有如下结论.

\begin{RiemannTheorem}\label{thm:length-minimizer-is-geodesic}
设 $(M,\langle\cdot,\cdot\rangle)$ 是 Riemann 流形, $p,q\in M$. 若 $c:I\to M$ 是连接 $p$ 到 $q$ 的一条分段可微曲线, 且对连接 $p$ 到 $q$ 的任意分段可微曲线 $\gamma:J\to M$, 都有 $l(c)\leq l(\gamma)$, 则 $c$ 是一条重参数化的测地线.
\end{RiemannTheorem}

\begin{RiemannProof}
为证明本定理, 需要以下定义.

\begin{RiemannDefinition}\label{def:totally-normal-neighborhood}
若存在 $\varepsilon>0$, 使得对所有 $p\in V$ 均有 $V\subset B_\varepsilon(p)$, 则称正规邻域 $V\subset M$ 为\textbf{全正规邻域}.
\end{RiemannDefinition}

现在证明全正规邻域总是存在. 为此, 回顾 $M$ 上的局部坐标 $(x^1,\ldots,x^n)$ 给出 $TM$ 上的局部坐标 $(x^1,\ldots,x^n,v^1,\ldots,v^n)$, 它们标记向量
\[
  v^1\frac{\partial}{\partial x^1}+\cdots+v^n\frac{\partial}{\partial x^n}.
\]
测地线方程
\[
  \ddot x^i+\sum_{j,k=1}^n\Gamma_{jk}^i\dot x^j\dot x^k=0\quad(i=1,\ldots,n),
\]
对应于一阶常微分方程组
\[
  \begin{cases}
    \dot x^i=v^i\\
    \dot v^i=-\sum_{j,k=1}^n\Gamma_{jk}^iv^jv^k
  \end{cases}
  \qquad(i=1,\ldots,n).
\]
这些方程定义了向量场 $X\in\mathfrak{X}(TM)$ 的局部流, 这个向量场在局部坐标下为
\[
  X=\sum_{i=1}^n v^i\frac{\partial}{\partial x^i}
  -\sum_{i,j,k=1}^n\Gamma_{jk}^iv^jv^k\frac{\partial}{\partial v^i},
\]
称为\textbf{测地流}. 如第~\ref{chap:differentiable-manifolds}~章所见, 对每个点 $v\in TM$, 存在开邻域 $W\subset TM$ 和包含 0 的开区间 $I\subset\mathbb{R}$, 使得 $X$ 的局部流 $F:W\times I\to TM$ 有良好定义. 特别地, 对每一点 $p\in M$ 可以选取包含 $p$ 的开邻域 $U$ 和 $\varepsilon>0$, 使得测地流在 $W\times I$ 上有良好定义, 其中
\[
  W=\{v\in T_qM\mid q\in U,\ \lVert v\rVert<\varepsilon\}.
\]
利用测地线的齐次性, 使 $\varepsilon$ 足够小, 就可以使区间 $I$ 任意大. 因此, 对足够小的 $\varepsilon$, 可以用 $G(v):=(q,\exp_q(v))$ 定义映射 $G:W\to M\times M$. 由于 $\exp_q(0)=q$, $(\mathrm{d}G)_0$ 在上述局部坐标下的矩阵表示为 $\bigl(\begin{smallmatrix}I&0\\I&I\end{smallmatrix}\bigr)$, 因而 $G$ 是局部微分同胚. 必要时缩小 $U$ 和 $\varepsilon$, 就可假设 $G$ 是到其像 $G(W)$ 的微分同胚, 而该像包含点 $(p,p)=G(0_p)$. 取 $p$ 的开邻域 $V$, 使 $V\times V\subset G(W)$, 则 $V$ 显然是全正规邻域: 对每个 $q\in V$, 有 $\{q\}\times\exp_q(B_\varepsilon(0))=G(W)\cap(\{q\}\times M)\supset\{q\}\times V$, 即 $\exp_q(B_\varepsilon(0))\supset V$.

注意, 给定全正规邻域 $V$ 中任意两点 $p,q$, 存在连接 $p$ 到 $q$ 的测地线 $c:I\to M$, 使得任意其他连接 $p$ 到 $q$ 的分段可微曲线 $\gamma:J\to M$ 都满足 $l(\gamma)\geq l(c)$ (且 $l(\gamma)=l(c)$ 当且仅当 $\gamma$ 是 $c$ 的重参数化). 现在, Theorem~\ref{thm:length-minimizer-is-geodesic} 的证明由以下观察直接得到: 若 $c:I\to M$ 是连接 $p$ 到 $q$ 的分段可微曲线, 且对连接 $p$ 到 $q$ 的任意曲线 $\gamma:J\to M$ 均有 $l(c)\leq l(\gamma)$, 则 $c$ 在它所经过的每个全正规邻域内必是一条重参数化的测地线.
\end{RiemannProof}

\begin{RiemannExercise}\label{ex:minimizing-properties-of-geodesics}
\begin{enumerate}[label=(\arabic*),ref=\arabic*,leftmargin=*]
\item\label{ex:unit-gradient-integral-curves} 设 $(M,g)$ 是 Riemann 流形, $f:M\to\mathbb{R}$ 是光滑函数. 若 $\lVert\operatorname{grad}(f)\rVert\equiv1$, 则 $\operatorname{grad}(f)$ 的积分曲线是测地线. 分别利用下列两种方法证明:
\begin{enumerate}[label=(\alph*),leftmargin=*]
\item 测地线的定义;
\item 测地线的最短性.
\end{enumerate}
\item\label{ex:metric-and-connection-in-normal-coordinates} 设 $M$ 是 Riemann 流形, $\nabla$ 是 $M$ 上的 Levi–Civita 联络. 给定 $p\in M$ 和 $T_pM$ 的一组基 $\{v_1,\ldots,v_n\}$, 考虑正规邻域的参数化 $\varphi:U\subset\mathbb{R}^n\to M$, 它由
\[
  \varphi(x^1,\ldots,x^n)=\exp_p(x^1v_1+\cdots+x^nv_n)
\]
给出 (局部坐标 $(x^1,\ldots,x^n)$ 称为\textbf{正规坐标}). 证明:
\begin{enumerate}[label=(\alph*),leftmargin=*]
\item 在这些坐标下, $\Gamma_{jk}^i(p)=0$ {\small (\textbf{提示:} 考虑测地线方程)};
\item 若 $\{v_1,\ldots,v_n\}$ 是标准正交基, 则 $g_{ij}(p)=\delta_{ij}$.
\end{enumerate}
\item\label{ex:bi-invariant-metrics-on-lie-groups} 设 $G$ 是配备\textbf{双不变 Riemann 度量}的 Lie 群 (即对所有 $g\in G$, $L_g$ 和 $R_g$ 都是等距映射), 并设 $i:G\to G$ 是由 $i(g)=g^{-1}$ 定义的微分同胚.
\begin{enumerate}[label=(\alph*),leftmargin=*]
\item 计算 $(\mathrm{d}i)_e$, 并证明对所有 $g\in G$, 有
\[
  (\mathrm{d}i)_g=(\mathrm{d}R_{g^{-1}})_e(\mathrm{d}i)_e(\mathrm{d}L_{g^{-1}})_g.
\]
由此推得 $i$ 是等距映射.
\item 设 $v\in\mathfrak{g}=T_eG$, $c_v$ 是满足 $c_v(0)=e$ 和 $\dot c_v(0)=v$ 的测地线. 证明若 $t$ 足够小, 则 $c_v(-t)=(c_v(t))^{-1}$. 由此推得 $c_v$ 在 $\mathbb{R}$ 上有定义, 且对所有 $t,s\in\mathbb{R}$ 满足 $c_v(t+s)=c_v(t)c_v(s)$.
{\small (\textbf{提示:} 回顾全正规邻域中的任意两点由该邻域内唯一一条测地线连接).}
\item \phantomsection\label{idx:bi-invariant-metric-geodesics}证明 $G$ 的测地线是左不变向量场的积分曲线, 并证明映射 $\exp$ (Lie 群的指数映射) 和 $\exp_e$ (单位元处的测地指数映射) 一致.
\item 设 $\nabla$ \phantomsection\label{idx:bi-invariant-metric-levi-civita}是双不变度量的 Levi-Civita 联络, $X,Y$ 是两个左不变向量场. 证明
\[
  \nabla_XY=\frac12[X,Y].
\]
\item 核实 Exercise~\ref{ex:riemannian-manifolds}(\ref{ex:left-invariant-metric}) 中的左不变度量实际上是双不变的.
\item 证明任意紧 Lie 群都容许双不变度量.
{\small (\textbf{提示:} 将一个左不变度量对所有右平移取平均).}
\end{enumerate}
\item\label{ex:isometries-send-geodesics-to-geodesics} 利用 Theorem~\ref{thm:length-minimizer-is-geodesic} 证明: 若 $f:M\to N$ 是等距映射且 $c:I\to M$ 是测地线, 则 $f\circ c:I\to N$ 也是测地线.
\item\label{ex:fixed-point-submanifold-of-isometry} 设 $f:M\to M$ 是等距映射, 其不动点集合是连通的 1 维子流形 $N\subset M$. 证明 $N$ 是一条测地线的像.
\item\label{ex:jacobi-fields-and-cut-locus} 设 $(M,\langle\cdot,\cdot\rangle)$ 是 Riemann 流形, 其测地线可以延拓到参数的所有取值. 设 $p\in M$.
\begin{enumerate}[label=(\alph*),leftmargin=*]
\item 设 $X$ 和 $Y_i$ 是以 $p$ 为中心的正规球上由 \eqref{eq:radial-field-in-normal-ball} 和 \eqref{eq:angular-field-in-normal-ball} 定义的向量场. 证明 $Y_i$ 满足 \textbf{Jacobi 方程}
\[
  \nabla_X\nabla_XY_i=R(X,Y_i)X,
\]
其中由
\[
  R(X,Y)Z=\nabla_X\nabla_YZ-\nabla_Y\nabla_XZ-\nabla_{[X,Y]}Z,
\]
定义的 $R:\mathfrak{X}(M)\times\mathfrak{X}(M)\times\mathfrak{X}(M)\to\mathfrak{X}(M)$ 称为\textbf{曲率算子} (参见第~\RGref{chap:curvature}~章).
{\small (\textbf{Remark:} 可以证明 $(R(X,Y)Z)_p$ 只依赖于 $X_p,Y_p,Z_p$).}
\item 考虑满足 $c(0)=p$ 的弧长参数化测地线 $c:\mathbb{R}\to M$. 若沿 $c$ 的向量场 $Y$ 满足沿 $c$ 的 Jacobi 方程
\[
  \frac{D^2Y}{\mathrm{d}t^2}=R(\dot c,Y)\dot c,
\]
则称它为 \textbf{Jacobi 场}. 证明 $Y$ 是满足 $Y(0)=0$ 的 Jacobi 场当且仅当
\[
  Y(t)=\left.\frac{\partial}{\partial s}\exp_p(tv(s))\right|_{s=0}
\]
其中 $v:(-\varepsilon,\varepsilon)\to T_pM$ 满足 $v(0)=\dot c(0)$.
\item 若点 $q\in M$ 是 $\exp_p$ 的临界值, 则称 $q$ \textbf{共轭于} $p$. 证明 $q$ 共轭于 $p$ 当且仅当存在一个非零 Jacobi 场 $Y$, 它沿着连接 $p=c(0)$ 和 $q=c(b)$ 的测地线 $c$ 且满足 $Y(0)=Y(b)=0$. 由此推知, 若 $q$ 共轭于 $p$, 则 $p$ 共轭于 $q$.
\item 若对所有 $X,Y\in\mathfrak{X}(M)$ 都有 $\langle R(X,Y)X,Y\rangle\geq0$, 则称流形 $M$ 具有\textbf{非正曲率}. 证明在这样的流形上没有两点共轭.
\item 给定一条由弧长参数化且满足 $c(0)=p$ 的测地线 $c:I\to M$, 设 $t_c$ 为使 $c$ 是连接 $p$ 到 $c(t)$ 的最短曲线的 $t$ 值所成集合的上确界 (因此 $t_c>0$). 定义 $p$ 的\textbf{割迹}为所有形如 $c(t_c)$ 且 $t_c<+\infty$ 的点所成的集合. 当 $M$ 分别为以下流形时, 求给定点 $p\in M$ 的割迹:
\begin{enumerate}[label=(\roman*),leftmargin=*]
\item 带平直方形度量的环面 $T^n$;
\item 带标准度量的球面 $S^n$;
\item 带标准度量的射影空间 $\mathbb{R}P^n$.
\end{enumerate}
在这些例子中验证: 割迹中的任一点或者共轭于 $p$, 或者由两条长度相同但像不同的测地线与 $p$ 相连.
{\small (\textbf{Remark:} 这是割迹的一般性质, 证明见 \RGcite{ch03-dC93} 或 \RGcite{ch03-GHL04}).}
\end{enumerate}
\end{enumerate}
\end{RiemannExercise}
